%% ===================================================================
%% Chapter 1. Introduction
%% Structure: Universe (broad context) -> Galaxy (research field and gap)
%%            -> Star (this thesis)
%% ===================================================================

\chapter{Introduction}

%% -------------------------------------------------------------------
%% Universe: why intercepting drones matters
%% -------------------------------------------------------------------
\section{Background: Drones as Both Tool and Threat}
\noindent
Small unmanned aerial vehicles (UAVs) have moved from research laboratories into almost every sector of society, including logistics, agriculture, infrastructure inspection, and public safety. The same properties that make them useful---low cost, high agility, and ease of operation---also make them a growing security concern. Drones have been used for unauthorized surveillance, for disrupting airports and critical infrastructure, and, increasingly, as low-cost weapons in armed conflicts \cite{chauhanNationsDefenseComprehensive2025}. As a consequence, counter-UAS (C-UAS) capabilities have become a priority for both civilian airspace management and national defense \cite{chauhanNationsDefenseComprehensive2025, palmasReinforcementLearningDecisionLevel2025}.

Among the available countermeasures, physical interception by another drone is attractive because it neutralizes the target without the collateral risks of jamming or kinetic weapons in populated areas \cite{lucotteInterceptionAdversarialDrones2026, vrbaAutonomousCaptureAgile2022}. Classical interception systems rely on guidance laws such as proportional navigation together with explicit target-state estimation and prediction \cite{pliskaSafeMidAirDrone2024, vrbaAutonomousCaptureAgile2022}. These approaches are well understood and can be deployed on real platforms, but they implicitly assume that the target follows a trajectory that can be extrapolated from its recent motion. A non-cooperative drone that actively reacts to its pursuer violates this assumption: it can dodge, change strategy, or exploit the dynamic limits of the interceptor, turning interception into a genuinely adversarial game.

In parallel, deep reinforcement learning (RL) has transformed what is possible in agile quadrotor flight. Policies trained in simulation now perform acrobatic maneuvers \cite{kaufmannDeepDroneAcrobatics2021}, beat human world champions in drone racing \cite{kaufmannChampionlevelDroneRacing2023a, bahnamMonoRaceWinningChampionLevel2026, boselloYourOwnProLevel2026}, and transfer zero-shot to real hardware when trained with low-level collective-thrust-and-body-rate (CTBR) commands \cite{kaufmannBenchmarkComparisonLearned2022, chenWhatMattersLearning2025a}. GPU-parallel simulators such as OmniDrones \cite{xuOmniDronesEfficientFlexible2023} make it feasible to collect billions of interaction steps for such policies. These advances suggest that an interceptor policy can be learned end-to-end, directly from interaction with a moving target, rather than being hand-designed.

%% -------------------------------------------------------------------
%% Galaxy: learning-based pursuit-evasion and its open problems
%% -------------------------------------------------------------------
\section{Learning-Based Pursuit--Evasion and Its Challenges}
\noindent
RL has been applied to target interception in a variety of forms, from missile guidance \cite{huInterceptGuidanceManeuvering2023} and fixed-wing interception in controlled airspace \cite{giralInterceptingUnauthorizedAerial2024} to vision-based drone interception \cite{darwishDroneNavigationTarget2023}, multi-target tracking \cite{xiaUAVMultiDynamicTarget2024}, and cooperative multi-agent interception \cite{fanSwitchingawareMultiagentDeep2023, zhangGameDronesMultiUAV2023, chenOnlinePlanningMultiUAV2025}. Many of these works, however, train the pursuer against targets with fixed, scripted, or randomly generated behavior. A pursuer trained this way learns to exploit the specific regularities of its training targets and is not guaranteed to remain effective against an evader that behaves strategically.

Modeling the evader as a learning agent turns the problem into a two-player zero-sum Markov game, which falls in the domain of competitive multi-agent reinforcement learning (MARL) \cite{zhangMultiAgentReinforcementLearning2021, liangReviewMultiAgentReinforcement2025, loweMultiAgentActorCriticMixed2017}. Training both sides simultaneously, however, introduces well-known difficulties: each agent's environment becomes non-stationary because its opponent keeps changing; agents overfit to the current opponent and forget how to handle earlier strategies (catastrophic forgetting); and the dynamics of mutual best responses can lead to strategy cycling rather than steady improvement \cite{liFightLadderBenchmarkCompetitive2024, ronceroLearnedControllersAgile2026}. Population-based self-play addresses these issues by training against a growing pool of past and diverse opponents instead of only the latest one. This idea underlies league training in AlphaStar \cite{vinyalsGrandmasterLevelStarCraft2019}, opponent-sampling schemes in autonomous air combat \cite{chenAutonomousDogfightDecisionMaking2025}, and, more recently, quadrotor pursuit--evasion and racing \cite{xiaoLearningMultipursuitEvasion2024, ronceroLearnedControllersAgile2026, tangAgilePEAutonomousUAV2026a, gavinAgileInterceptionFlying2026, gavinInterceptingAgileTarget2026, gelesSuperhumanSafeAgile2026}. In these works, the choice of \emph{which} opponent to train against---uniformly, by prioritized fictitious self-play (PFSP), or by a Hedge-style sampler over win rates---has a decisive effect on the robustness of the final policy \cite{chenAutonomousDogfightDecisionMaking2025, ronceroLearnedControllersAgile2026, gavinInterceptingAgileTarget2026}.

Two gaps remain. First, existing opponent samplers rank opponents almost exclusively by win rate. In physical pursuit--evasion, however, win rate is a coarse signal: an evader can be ``hard'' because it escapes cleanly, because it drags the capture out over a long time, or because it forces the pursuer into aggressive, crash-prone flight. Conversely, an evader can appear easy simply because it destroys itself by flying into the ground or a wall. A sampler that only sees win rate cannot distinguish these cases and may focus training on the wrong opponents.

Second, even a well-trained population-based policy is \emph{frozen} at deployment. It encodes a response to the mixture of opponents it saw during training, but it has no explicit mechanism to identify and adapt to the particular evader it is currently facing, especially one whose strategy was never seen in training or that changes strategy mid-episode. Opponent modeling \cite{tianMultiagentActorcriticTime2023} and target-prediction networks \cite{zhangGameDronesMultiUAV2023, chenOnlinePlanningMultiUAV2025} supply the policy with an estimate of the opponent's future behavior, but they are typically trained offline and not updated online. Rapid online adaptation has been studied extensively for changing \emph{dynamics}, through residual learning, differentiable simulation, and latent-condition inference \cite{panLearningFlyRapid2026, xingContinualRobotPolicy2026, pertinNSTRLResidualReinforcement2026, greenBridgingPerformanceGeneralization2026}, yet these methods generally rely on gradient-based updates or privileged information, which are costly on resource-constrained flight hardware. Reservoir computing offers a lightweight alternative: an echo state network (ESN) with a fixed random reservoir and a linear readout updated by recursive least squares (RLS) can adapt within a few control steps without backpropagation or pretraining \cite{yoshimuraOnlineAdaptiveReinforcement2026}. To the best of our knowledge, such online adaptation has not been combined with population-based self-play for adversarial aerial pursuit.

%% -------------------------------------------------------------------
%% Star: this thesis
%% -------------------------------------------------------------------
\section{This Thesis: Adaptive Pursuit via Population-Based Self-Play and an ESN Module}
\noindent
This thesis develops an adaptive pursuit policy for one-versus-one quadrotor pursuit--evasion that addresses both gaps. The pursuer and the evader are both Crazyflie quadrotors controlled by PPO policies \cite{yuSurprisingEffectivenessPPO} that output CTBR commands tracked by an on-board rate controller, and they are trained in a GPU-parallel Isaac Sim environment built on OmniDrones \cite{xuOmniDronesEfficientFlexible2023} inside a bounded arena. The approach has two components.

The first component is a generational, population-based self-play framework. Starting from scripted heuristic opponents (hover, circular orbit, pursue, and flee), each generation trains one side against frozen opponents drawn from the other side's population, evaluates the new policy against the entire pool, and adds it back. Instead of sampling opponents by win rate alone, we score each opponent with a temperature-scaled softmax over several normalized \emph{hardness} metrics---capture rate, time to capture, motor effort, and the training agent's own crash and out-of-bounds rate---so that training concentrates on opponents that are difficult in different ways. A reward design with symmetric terminal outcomes, partial credit for opponent-induced failures, and body-rate command penalties prevents degenerate strategies such as waiting for the opponent to crash and holding saturated rate commands that would tumble a real vehicle.

The second component is an ESN-based online adaptation module attached to the pursuer. A fixed random reservoir encodes the pursuer's observation history, and a linear readout predicts the evader's next relative position and velocity in the pursuer's body frame. The readout is updated online by RLS at every control step, and its prediction is appended to the policy input. This gives the pursuer a continually updated, opponent-specific forecast at the cost of a matrix update per step, without gradient computation at test time. We study three readout variants: a shared readout fitted centrally during training and frozen at deployment, a per-episode readout adapted from scratch, and a shared readout used as the initialization for per-episode adaptation.

The main contributions of this thesis are summarized as follows:
\begin{itemize}
    \item A population-based self-play framework for agile quadrotor pursuit--evasion with CTBR control, including heuristic population seeding, a reward structure that removes degenerate strategies, and a multi-metric, temperature-controlled opponent sampler that extends win-rate-based sampling.
    \item The integration of an echo state network with online RLS readout into the pursuer policy as a lightweight, gradient-free opponent-prediction and adaptation module, together with shared, per-episode, and hybrid readout schemes.
    \item An evaluation protocol based on cross-generation win-rate matrices and held-out evader strategies, including strategy switches within an episode, to assess both the robustness of the population and the adaptation ability of the pursuer.
\end{itemize}

\section{Thesis Organization}
\noindent
The remainder of this thesis is organized as follows. Chapter~2 reviews related work on drone interception, competitive multi-agent reinforcement learning and self-play, and online adaptation. Chapter~3 formulates the pursuit--evasion problem and presents the proposed self-play framework and ESN adaptation module. Chapter~4 describes the experimental setup and discusses the results. Chapter~5 concludes the thesis and outlines directions for future work.
